\section{Edition 06 geometric suppliers and attained rigidity}
\label{sec:ed06}
These open formal milestones follow the master entry (E), fixed-type
attainment (X) and rigidity (R) plan. Use the oriented smooth carrier and
physical measures of Section~\ref{sec:ed05}. Thus $T=4\pi^2$,
$e=E_+=E_-$ on the closed smooth homotopy four-sphere $M$, and
$I=I^+(M)$. Write $D=\int|\operatorname{Ric}_0|^2\,dV$.
The assessed attainment argument in Investigation 06
supplies the scientific starting point for R. The formal attainment
and recognition-gap targets remain open in
Theorem~\ref{thm:ed05-bour-attainment} and
Proposition~\ref{prop:ed05-uniform-gap}.

\subsection{Entry and independent marked reduction}

\begin{openproblem}[Positive entry on a marked handlebody boundary]
\label{prob:ed06-entry-construction}
\uses{prop:ed04-positive-class,prop:ed04-round-face-control}
Route 11. Fix one actual balanced-presentation five-dimensional
handlebody with handles of indices zero, one and two whose oriented
boundary is a smooth homotopy four-sphere. Specify the attaching circles,
framings, retained exterior and a compatible new cap. The tested product
circle cap already costs at least $64\pi^2/3$ in the sum of the signs.
On the resulting closed metric prove $E_++E_-<16\pi^2$ by separately
accounting for both signs on every piece. Prove that for some $\lambda>0$
the conformal-Laplacian quadratic form dominates
$\lambda\int u^2\,dV_g$ for every smooth test $u$. Record the marked
smooth identification for this construction. An isometric
double comparison requires the actual boundary isometry and estimates
on the same physical half metric as in Problem~\ref{ed04:entry}.
On the marked Miller--Schupp boundary, the regular and long-cuff full-width trials still need a positive closed sign and the two signed costs. Coverage of additional smooth types is the separate target
in Theorem~\ref{thm:ed05-coverage}.
\end{openproblem}

\begin{openproblem}[Native reduction and its first uncovered cases]
\label{prob:ed06-marked-reduction}
\uses{prop:ed45-relative-pair,prop:ed46-reverse-core}
Route 09. Formalize the assessed Investigation 05 criterion on a smooth
four-ball with a split zero-framed lower unknot $A$, an integral pair of
native attachments from $Y_1$ to the immediate $Y_3$, and the actual
rank-three terminal upper body and closing, with
$Y_r=\#_r(S^1\times S^2)$. If either single-attachment cut of the pair is
reducible, the corresponding stationary knot localizes in a ball.
Use the stationary filling and framed slides to localize both tubes,
obtain the local surgery $Y_2$, put $A$ last and cancel with the actual
upper body. The assessed conclusion is a strict same-type
$(0,3,3)\to(0,2,2)$ reduction in indices one through three, with no
stabilization or hypothesis on the initial intersection count. Carry
the forward and inverse maps of the actual closing. With a residual
continuation, account for its transport and inverse before asserting
any decrease. In a least-three presentation of this shape the two
native second cuts are irreducible, have first Betti number two and
admit no nonzero-degree map from a closed oriented positive-scalar
three-manifold. Seek a marked slide, new order or level for one such
input, or recognize one actual localized rank-two closing. Route 10
owns the physical deformation on a supplied qualifying cut. The product-doubled tangle class now has a scientific recognition with three cancellations and no births, including a control with both cuts irreducible. On the native GST outlet, up to three slides do not yield a first unstabilized cancellation in the assessed family.
\end{openproblem}

\subsection{Physical suppliers for fixed-type attainment}

\begin{proposition}[First variation on the original metric]
\label{prob:ed06-original-variation}
\uses{prop:ed50-gradient,prop:ed50-commuted}
\lean{SmoothPoincare.Bour.originalGradient_pairing, SmoothPoincare.Bour.bourCovariantGradient_apply_commuted, SmoothPoincare.Bour.bourOriginalVelocity_eq_negTwoGradient_iff_components, SmoothPoincare.Bour.originalGradient_regularSlab}
\leanok
On a smooth closed oriented four-dimensional regular metric family with
constant Euler integral, the derivative of $\mathcal F^p$ pairs the
symmetric tensor $G^p$ of Proposition~\ref{prop:ed50-commuted} with
the actual metric velocity. Its explicit component evolution on every tangent
pair is equivalent to the physical equation $K=-2G^p$. A regular
solution of the negative twice-gradient equation uses this same tensor
for dissipation and volume preservation. Local conformal Euler constancy
also gives $\operatorname{tr}G^p=-p\Delta R/4$.
\end{proposition}
\begin{proof}
\leanok
Apply the two adjoints to the physical Ricci velocity, include both
inverse-metric and volume derivatives, then commute the crossed Ricci
derivatives. Test on the actual scale, conformal and solution paths.
\end{proof}

\begin{openproblem}[Harmonic radius and uniform stationary interior estimates]
\label{prob:ed06-stationary-interior}
\uses{prop:ed51-morrey,prop:ed47-mixed,prop:ed48-trace}
Route 13. Use the genuine constant-scalar stationary equation for each
mixed return $q_j$, with $p_j=\alpha_j$ for a closed regularized return
and $p_j=0$ for an exact Yamabe return. On the original and
curvature-normalized scales $h_j=a_jq_j$, $a_j\ge1$, start with a common
ordered four-slot curvature bound and the measured volume and injectivity
estimates. Supply the whole-space scalar $L^6$ Laplacian estimate in the pinned library and construct a common interior harmonic radius with controlled
$C^{1,\gamma}$ coefficients. Use the scalar diagonal Ricci principal
part in these coordinates and the $\tfrac12\nabla^*\nabla$ principal
part of the trace-free Ricci equation to obtain coupled interior
$W^{2,b}$ estimates for $b>4$ and then estimates of every fixed order
on smaller balls. The coupled physical estimate at exponent six is Proposition~\ref{prop:ed51-coupled}. Control $s_j=R(q_j)/a_j$ and the coefficient
$(2-p_j)s_j/12$ uniformly, with no division by $p_j$.
\end{openproblem}

\begin{openproblem}[Nonflat ALE end geometry and marked return]
\label{prob:ed06-ale-return}
\uses{prop:ed51-pfaffian,prop:ed49-exterior,prop:ed47-mark,prop:ed47-ends,prop:ed48-transfer,prop:ed48-cap}
Route 13. Treat the fixed-parameter original extraction and the mixed
stationary extraction at their respective scales. Given a complete
nonflat scalar-flat Bach-flat pointed limit with finite full curvature,
positive Yamabe test bounds and the source chiral and Ricci budgets,
prove ALE decay and the limiting Gauss--Bonnet boundary terms including
ambient curvature. Exclude nontrivial link groups before applying
$D_X=2(E_+(X)+E_-(X))+32\pi^2(k-1)$ to exclude $k\ge2$ by the
strict Ricci ceiling. On one connected truncated core with all end
buffers, obtain a coherent oriented embedding, actual collar-compatible
cap data and forward and inverse smooth comparisons. Construct a
positive compactification on the same smooth type with
$I\le e_X\le\liminf_i e_i$ and $D_X=4e_X$. The separate
start-deficit test of Proposition~\ref{prop:ed48-deficit} applies to
a second marked return with the specified tail hypothesis.
\end{openproblem}

\subsection{Rigidity on the attained exact stratum}

\begin{openproblem}[Least Ricci selector at the exact attained value]
\label{prob:ed06-least-ricci}
\uses{prop:ed49-selector,thm:ed05-bour-attainment,prop:ed48-transfer}
Route 07. Fix one entered nonstandard smooth type with attained
$I=a_M$ and $T+\varepsilon_*<I\le 2T-\delta$ for some fixed $\delta>0$,
where $a_M$ is the first positive critical value and $\varepsilon_*$ is
the recognition gap. Formalize
the compactness supplier for unit-volume positive Yamabe representatives
at the exact value $a_M$, modulo oriented diffeomorphisms, with the
fixed-band marked-neck and same-type strict-drop hypotheses. Pass the
Yamabe inequality and equality for constant tests to the limit. Then
$D=4I-(Y_{\rm rd}^2-Y^2)/12$, where $Y_{\rm rd}$ is the spherical
Yamabe constant. Select the minimum of $Y^2$ and hence of $D$.
Strict spherical comparison at the maximum of $Y^2$ bounds the deficit
away from zero on this exact stratum. Relate that bound to near-minimizing
starts only under a separate recovery hypothesis. The post-attainment value derivative is computed in Proposition~\ref{prop:ed51-slope}.
\end{openproblem}

\begin{openproblem}[Polarized Weyl Hessian for a non-Einstein minimizer]
\label{prob:ed06-polarized-hessian}
\uses{prob:ed06-least-ricci,prob:ed06-original-variation,prop:ed47-variation}
Route 07. On the least-$D$ candidate with $D>0$ let
$Z=\operatorname{Ric}_0$ and
$P=(Z^\sharp\circ Z^\sharp)^\flat-\tfrac14|Z|^2g$ in its
constant-scalar Yamabe gauge. Compute the polarized physical Hessian
$H=D^2E_+$ on $Z,P$ modulo conformal and diffeomorphism directions.
Identify a geometric subclass where the span is independent modulo
gauge and the Hessian matrix has negative determinant, or where
$H(Z,Z)=0$ and $H(Z,P)\ne0$. Realize the resulting negative direction
by an admissible metric path and a strictly lower positive same-type
competitor. Record the dependent and nonnegative cases that remain.
\end{openproblem}

\begin{openproblem}[A transverse-traceless spectral window in the Einstein case]
\label{prob:ed06-tt-window}
\uses{prop:ed49-round,prob:ed06-least-ricci,prop:ed47-variation}
Route 07. For $D=0$ the selected positive Yamabe metric is Einstein.
At scalar $s>0$, the physical Weyl Hessian on transverse-traceless
tensors $h$ is $\tfrac18\langle h,(-\Delta_L-s/2)(-\Delta_L-s/3)h\rangle_{L^2}$,
where $\Delta_L$ has leading term $\nabla^i\nabla_i$ in the master
convention. For a specified Einstein subclass, force an actual smooth
transverse-traceless eigentensor of $-\Delta_L$ with eigenvalue in
$(s/3,s/2)$ and realize its admissible negative variation, or give
smooth recognition on that subclass. Endpoint eigentensors have zero
Hessian.
\end{openproblem}

\subsection{Checked results on the physical metric and the marked exterior}

\begin{proposition}[The two spatial adjoints of the connection variation]
\label{prop:ed49-adjoint}
\uses{prop:ed48-slab,prop:ed48-trace}
\lean{SmoothPoincare.Bour.bourRicciDivergence_adjoint, SmoothPoincare.Bour.bourWeightedRicciKoszul_firstAdjoint, SmoothPoincare.Bour.bourWeightedRicciKoszul_secondAdjoint, SmoothPoincare.Bour.bourKoszulFlux_native, SmoothPoincare.Bour.bourClosedTensorPairing_native}
\leanok
At a smooth regular time on a closed physical metric path, write $K=\dot g$,
$A=\operatorname{Ric}+(p-4)Rg/12$, and
$C_{kij}=\tfrac12(\nabla_i K_{jk}+\nabla_j K_{ik}-\nabla_k K_{ij})$.
Then $C$ is the lowered variation of the Levi--Civita connection.
The physical $L^2$ pairing of $\nabla^k C_{kij}$ with $A$ is
$-\langle C,\nabla A\rangle_{L^2}$, with the derivative in the first slot.
Moving the remaining derivative gives a pairing with $K$, using the inverse
slot permutations and coefficient $-1/2$ on the rough Laplacian of $A$.
\end{proposition}
\begin{proof}
\leanok
Apply covariant divergence integration by parts twice, using invariance of the
metric pairing under simultaneous index permutation. The differentiated Koszul
formula and the native smooth velocity identify the connection term on the path.
\end{proof}

\begin{proposition}[The time derivative of the original Ricci tensor in a fixed chart]
\label{prop:ed49-ricci-time}
\uses{prop:ed48-slab}
\lean{SmoothPoincare.Bour.bourChartChristoffel_spatial_hasDerivAt, SmoothPoincare.Bour.bourChartRicci_hasDerivAt, SmoothPoincare.Bour.bourMetricRicciAt_chartBasis_hasDerivAt}
\leanok
On a smooth regular metric family in a fixed chart, time differentiation
commutes with each spatial derivative of the actual Christoffel coefficients.
The derivative of each physical Ricci component on the fixed chart basis
is obtained by taking the two spatial derivatives of the Christoffel speed
in their original order and differentiating both Christoffel products.
\end{proposition}
\begin{proof}
\leanok
Use joint smoothness in time and space on the chart interior to exchange
the derivatives. Differentiate the two contracted products by the product
rule. The chart-basis curvature identity identifies the result with the
corresponding component of the original Ricci tensor.
\end{proof}

\begin{proposition}[Connection estimates in a supplied harmonic chart]
\label{prop:ed49-interior}
\uses{prop:ed48-harmonic,prop:ed48-normal}
\lean{SmoothPoincare.bourChartConnectionDerivative_bound, SmoothPoincare.bourChartConnection_trace_eq_zero, SmoothPoincare.bourChartRoughLaplacian_eq_secondCovariantTrace}
\leanok
In the four-dimensional metric chart with inverse entries bounded by $A$,
first derivatives by $L$ and second derivatives by $D$,
each first derivative of a Christoffel coefficient is bounded by
$96A^2 L^2+6AD$. At a harmonic point the traced coordinate-frame
connection vanishes. The raw tensor connection Laplacian on a smooth
compactly supported tensor reduces there to the metric trace of the
two iterated coordinate covariant derivatives.
\end{proposition}
\begin{proof}
\leanok
Differentiate the inverse-metric identity and the Christoffel formula.
Contract the Laplacians of the coordinate functions, then use the
linearity of the tensor connection in its directional input.
\end{proof}

\begin{proposition}[Euler transgression on a supplied nonflat end]
\label{prop:ed49-end}
\uses{prop:ed48-cap,prop:ed48-transfer}
\lean{SmoothPoincare.bourEulerBoundaryDensity_withArea_error, SmoothPoincare.bourALELevelBoundary_tendsto, SmoothPoincare.bourExhaustedEulerIdentity, SmoothPoincare.bourOneEuclideanEnd_of_boundaryLimits}
\leanok
Use the outward shape and the curvature matrix on the three
complementary oriented tangent two-planes of a normalized link.
If the shape, curvature and area ratio are within $\varepsilon\le1$
entrywise of $\mathrm{Id},0,1$, the area-weighted boundary density
differs from $4$ by at most $412\varepsilon$.
With normalized error tending to zero, integrability and the degree-$m$
round covering measure, its integral tends to $8\pi^2/m$.
For a supplied scalar-flat compact exhaustion with the finite Euler
formula, spherical links, $\chi=2-k$, and $0\le D<32\pi^2$, the
nonempty set of ends has cardinality one.
\end{proposition}
\begin{proof}
\leanok
Expand the cubic determinant and bound the mixed shape--curvature contraction.
Integrate the error on the fixed quotient link and pass all $k$ boundary
terms through the limit. The resulting identity is
$D=2(E_++E_-)+32\pi^2(k-1)$.
\end{proof}

\begin{proposition}[Buffered disk extension on a recognized sphere]
\label{prop:ed49-exterior}
\uses{prop:ed48-cap}
\lean{SmoothPoincare.bourRecognizedCap_commonChart, SmoothPoincare.bourMarkedExterior_of_recognizedSphere}
\leanok
Let the capped source exterior and target double be orientedly recognized
as standard four-spheres. Suppose their cap coordinates extend over a closed
radius-$R$ disk, the seam has radius $0<r<R$, and the two parametrizations
carry the same center orientation. If their composition gives the prescribed
map on the buffered disk, an oriented global diffeomorphism agrees with it
and with its inverse on the two open buffers and identifies the complements
of the open radius-$r$ disks.
\end{proposition}
\begin{proof}
\leanok
Choose an exterior point for each full buffer. On the standard sphere,
move the first point to the second by an oriented map and align the cap
centers by a further motion fixing that exterior point. Stereographic
projection now gives a common chart on the entire buffers. At their common
center the transported frames have equal orientation. Apply the supported
cap extension in this chart and pull it back to the given double.
\end{proof}

\begin{proposition}[Least Ricci energy from first-value compactness]
\label{prop:ed49-selector}
\uses{prop:ed48-transfer}
\lean{SmoothPoincare.attainedYamabe_limit_of_metricConvergence, SmoothPoincare.exists_attainedYamabe_leastRicci}
\leanok
Supply positive subcritical attainment at $0<I\le8\pi^2-\delta$ on the
fixed type, conformal normalization and Bach covariance, exact first-value
compactness modulo oriented pullback, the closed Euler value and the sharp
spherical comparison with equality implying zero chiral energy.
Then $I=a_M$. One unit-volume exact Yamabe representative minimizes
both $Y^2$ and $D$. Another attains the maximum $Y^2$ and its strictly
positive squared spherical deficit bounds the deficit on the exact stratum.
\end{proposition}
\begin{proof}
\leanok
Pass each fixed smooth conformal test through smooth metric convergence and
use the constant test for equality. Apply compactness to a minimizing
and a maximizing sequence separately and use the exact closed Euler identity.
\end{proof}

\begin{proposition}[The round transverse-traceless window and a conditional descent]
\label{prop:ed49-round}
\uses{prop:ed47-variation}
\lean{SmoothPoincare.einsteinWeylTTFactor_neg_iff, SmoothPoincare.standardRoundFour_curvatureAction_traceFree, SmoothPoincare.standardRoundFour_negLichnerowicz_eigenvalue_ge_eight, SmoothPoincare.standardRoundFour_TTMode_weylFactor_ge_one, SmoothPoincare.bourStrictTTWindowExcludesAttainment}
\leanok
On the unit round sphere the trace-free operator $8\mathrm{Id}-\operatorname{tr}\nabla^2$
has eigenvalues $\lambda\ge8$ on nonzero smooth modes. The coefficient
$(s/2-\lambda)(s/3-\lambda)/8$ at $s=12$ is at least one.
On any actual admissible path realizing this signed Hessian on a nonzero
tensor, a strict TT mode $s/3<\lambda<s/2$ excludes attainment.
\end{proposition}
\begin{proof}
\leanok
Compute the round curvature action. Pair the shifted rough operator with
the eigentensor and use the covariant Green identity. Factor the signed
coefficient and use the physical two-sided variation test at a strict sign.
\end{proof}

\begin{proposition}[Regularized Ricci margin from an attained comparator]
\label{prop:ed49-margin}
\uses{prop:ed47-mixed}
\lean{SmoothPoincare.Bour.bourOriginal_attainedComparator_margin, SmoothPoincare.Bour.bourNoMarkedBubble_of_ricciMargin}
\leanok
Fix $0<p<1$ and a comparator of energy $I$ on the original type.
On a supplied regularized trajectory suppose
$F^p(g_0)\le F^p(g_*)+\epsilon<(1-p)16\pi^2$ and
$4\epsilon\le p(4I-D_*)$. Every original smooth time has positive Yamabe
sign, $D(t)\le(4I+D_*)/2$, and $S(t)\le S_*+24\epsilon/p$.
Any sequence of original-type metrics with eventual $D\le d<4I$ excludes
a positive marked return satisfying $4e_X\le\liminf D$.
\end{proposition}
\begin{proof}
\leanok
The strict pinching margin keeps the full trajectory in the positive class.
Insert $e(t)\ge I$ in the two physical formulas for $F^p$.
A marked positive metric on the original type has $e_X\ge I$.
\end{proof}

\subsection{The physical gradient and both extractions}

\begin{proposition}[Ricci velocity on the original path]
\label{prop:ed50-palatini}
\uses{prop:ed49-adjoint,prop:ed49-ricci-time}
\lean{SmoothPoincare.Bour.bourMetricRicciAt_chartKoszul_hasDerivAt, SmoothPoincare.Bour.bourMetricRicciAt_hasDerivAt, SmoothPoincare.Bour.bourKoszulFlux_native_trace, SmoothPoincare.Bour.bourKoszulTraceForm_native}
\leanok
Let $K=\dot g$ on an original regular slice and
$C_{aik}=(\nabla_iK_{ka}+\nabla_kK_{ia}-\nabla_aK_{ik})/2$.
In the fixed chart axes and intrinsically,
$\dot{\mathrm{Ric}}_{ik}=\nabla^a C_{aik}-\nabla_k\tau_i$
with $\tau=\tfrac12 d\operatorname{tr}_gK$.
\end{proposition}
\begin{proof}
\leanok
Differentiate the chart connection and lower its output in the same chart.
Metric compatibility combines the differentiated inverse and connection
terms into the covariant derivative of the lowered flux. Trace output
and last slot to obtain $\tau$, then use a basis to recover the tensor.
\end{proof}

\begin{proposition}[A symmetric gradient and its regular time slab]
\label{prop:ed50-gradient}
\uses{prop:ed50-palatini,prop:ed49-adjoint,prop:ed48-slab,prop:ed48-trace}
\lean{SmoothPoincare.Bour.bourEnergy_hasDerivAt_contractions, SmoothPoincare.Bour.bourWeightedPalatini_secondAdjoint, SmoothPoincare.Bour.originalGradient_pairing, SmoothPoincare.Bour.bourCovariantGradient_trace, SmoothPoincare.Bour.originalGradient_regularSlab}
\leanok
Write $b=(p-4)/12$, $A=\mathrm{Ric}+bRg$ and
$B=(|\mathrm{Ric}|^2+bR^2)/2$.
On a closed regular signed metric family of constant Euler integral,
the covariant $L^2$ gradient is the symmetric part of
\[
 \nabla^a\nabla_i A_{aj}-\tfrac12\nabla^a\nabla_a A_{ij}
 -\tfrac12(\operatorname{div}\operatorname{div}A)g_{ij}
 -(\mathrm{Ric}\circ\mathrm{Ric})_{ij}-bR\,\mathrm{Ric}_{ij}+\tfrac12 B g_{ij}.
\]
Its integrated trace vanishes. Under local constancy of the Euler integral
on conformal tests its trace at each point is $-p\,\Delta R/4$ in the
convention $\Delta=\operatorname{tr}\nabla^2$. On any regular solution
with $\dot g=-2G^p$, volume is constant, and the integrated squared
gradient is half the energy drop.
\end{proposition}
\begin{proof}
\leanok
Differentiate the two inverse metrics in the Ricci square and the
inverse metric in the scalar trace. Include the volume derivative in
the same integral. Transfer both derivatives of each Palatini term.
The symmetric projection pairs with the original velocity. Use actual
scale and conformal paths for the trace tests and integrate along the given solution.
\end{proof}

\begin{proposition}[Physical capture by a common harmonic coordinate ball]
\label{prop:ed50-chart}
\uses{prop:ed48-normal,prop:ed48-harmonic,prop:ed49-interior}
\lean{SmoothPoincare.bourUniformHarmonicBalls_of_yamabe, SmoothPoincare.BourHarmonicChart.exists_buffered, SmoothPoincare.BourHarmonicChart.transport_to_original, SmoothPoincare.BourHarmonicChart.exists_model_metric, SmoothPoincare.bourUniformHarmonicMetricTrace_bound}
\leanok
Under the quantitative Peters harmonic-radius input at a common positive
injectivity radius and ordered curvature ceiling, the normalized metrics
$h=ag$ on the stationary sequence have one radius $r$ and one ellipticity
constant $L$. The chart is defined on $B(0,4r)$ and its closed radius-$r$
image contains the physical $h$-ball of radius $r/L$. The same set
contains the $g$-ball of radius $r/(L\sqrt a)$. An ambient model metric
agrees with the actual pullback on $B(0,3r)$, where the model's fixed
coordinates are harmonic. Its metric entries there satisfy a uniform
pointwise bound on $|h^{ab}\partial_a\partial_b h_{ik}|$ from the actual
ordered curvature ceiling, the first jets and ellipticity.
\end{proposition}
\begin{proof}
\leanok
Apply the physical noncollapse and injectivity estimates before choosing the
Peters chart. Two length comparisons and first exit give the nested balls.
Extend the pullback on the compact inner ball and use Laplacian naturality.
Pull Ricci back on the same ball and insert the curvature and first-jet
bounds in the harmonic metric equation.
\end{proof}

\begin{proposition}[Measured finite Euler correction on an end link]
\label{prop:ed50-link}
\uses{prop:ed49-end}
\lean{SmoothPoincare.bourUnitEuclideanLevel_shape_error, SmoothPoincare.bourImmersedLinkAreaDensity_error, SmoothPoincare.bourImmersedLinkEulerIntegral_error, SmoothPoincare.bourEulerBoundaryDensity_eq_integral_connectionPath}
\leanok
On the normalized unit sphere use $u=|x|^2/2$ and the outward normal.
Under the stated bilinear metric, connection difference, and ordered
ambient curvature bounds of size $\epsilon$, the shape error is at most
$20\epsilon$ and the induced area ratio error at most $3\epsilon$.
On a compact immersed link with the reference unit shape and
outward-adapted orthonormal frames, the physical boundary integral differs
from four times the reference area by at most
$412(20\epsilon)$ times that area when $20\epsilon\le1$.
The scalar connection-path integral has value
$2\sum S_{ab}C_{ab}+4\det S$.
\end{proposition}
\begin{proof}
\leanok
Compare the two Hessians of $u$, bound the length of its gradient and
compare the two induced volume measures. The norm of the ordered curvature
controls each complementary component. Integrate the E49 density estimate
on this fixed link. The tangential quadratic term in the connection
interpolation has weight $1-t^2$.
\end{proof}

\begin{proposition}[The same post-attainment Ricci margin at both extractions]
\label{prop:ed50-second}
\uses{prop:ed49-selector,prop:ed49-margin,prop:ed48-transfer,prop:ed47-return}
\lean{SmoothPoincare.Bour.exists_bourLeastRicci_uniformRegularizedStarts, SmoothPoincare.Bour.bourFirstClosedExtraction_bounds, SmoothPoincare.Bour.bourSecondMarkedExtraction_of_canonicalRescaling_impossible, SmoothPoincare.Bour.bourRegularizedMinimum_of_secondClosed, SmoothPoincare.Bour.bourRegularizedMinima_limit_leastRicci}
\leanok
Given scientific attainment, first-value compactness, the Yamabe and
sharp spherical comparison inputs on a type with $0<I\le8\pi^2-\delta$,
choose one exact least-$D$ metric $g_*$. The same
$d=(4I+D_*)/2<4I$ works for arbitrarily accurate near-minimizers at
every sufficiently small positive $p$. At each fixed $p$, physical closed
endpoints inherit the bound. A second marked positive return on the
original type with the one-end Ricci identity contradicts it.
A supplied whole closed second limit with the stated test and integral
convergences attains the regularized infimum. Along actual global
minima as $p\downarrow0$, a supplied positive whole Yamabe and Bach
limit attains the least $D$ and the chiral excess is $o(p)$.
\end{proposition}
\begin{proof}
\leanok
Compare both physical formulas for the energy with the same exact metric.
Transfer the bound in the original gauges along each first closed extraction.
For a second pointed limit the ordinary Ricci lower bound gives
$4I\le4e_X=D_X\le d$. For whole convergence use both sides of the
variational infimum bound. Finally compare each global minimizer with $g_*$
and pass its Ricci integral to the whole limit.
\end{proof}

\begin{proposition}[Vanishing dissipation on terminal backward windows]
\label{prop:ed50-terminal}
\uses{prop:ed50-gradient,prop:ed47-motion}
\lean{SmoothPoincare.Bour.originalGradient_energy_tendsto_terminal, SmoothPoincare.Bour.originalGradient_dissipation_tendsto_terminal, SmoothPoincare.Bour.originalGradient_parabolicDissipation_tendsto_terminal}
\leanok
For $0\le p\le1$, assume the original equation, Euler constancy, and an
order-connected regular time interval. The energy tends to its infimum
at the right end of this interval. If both endpoints of original slabs
are cofinal there, the dissipation on those slabs tends to zero. The same
holds on the corresponding fixed backward parabolic windows.
\end{proposition}
\begin{proof}
\leanok
The energy is antitone and the chiral--Ricci formula bounds it below by zero.
Take the difference of the two cofinal energy limits and apply the exact
rescaling identity to each window.
\end{proof}


\begin{proposition}[Curvature-commuted gradient on a general slice]
\label{prop:ed50-commuted}
\uses{prop:ed50-gradient}
\lean{SmoothPoincare.Bour.bourRicciWeight_divergence, SmoothPoincare.Bour.bourRicciWeight_doubleDivergence, SmoothPoincare.Bour.bourCrossedRicci_eq_ricciSquare_sub_curvature, SmoothPoincare.Bour.bourRicciFluxAdjoint_apply_commuted, SmoothPoincare.Bour.bourCovariantGradient_apply_commuted, SmoothPoincare.Bour.bourCovariantGradient_apply_csc}
\leanok
For $b=(p-4)/12$, $B=(|\mathrm{Ric}|^2+bR^2)/2$ and
$\Delta=\operatorname{tr}\nabla^2$, the same physical tensor satisfies
\[
\begin{aligned}
G^p={}&-\tfrac12\Delta\mathrm{Ric}-\mathrm{Rm}^{\circ}(\mathrm{Ric})
       +\tfrac{p+2}{12}\nabla^2R\\
     &-\tfrac{p-1}{12}(\Delta R)g-bR\,\mathrm{Ric}+\tfrac12Bg .
\end{aligned}
\]
For constant scalar curvature the two scalar derivative terms vanish.
\end{proposition}
\begin{proof}
\leanok
The contracted Bianchi identity gives
$\operatorname{div}(\mathrm{Ric}+bRg)=(p+2)dR/12$.
In an orthonormal basis the crossed curvature action is
$\mathrm{Ric}^2-\mathrm{Rm}^{\circ}(\mathrm{Ric})$.
Its first term cancels the moving-inverse correction. Use metric
compatibility to differentiate $Rg$ twice and trace its Hessian.
\end{proof}

\subsection{Harmonic interiors and terminal spacetime limits}

\begin{proposition}[Metric estimate on the physical pullback]
\label{prop:ed51-metric}
\uses{prop:ed50-chart}
\lean{SmoothPoincare.Bour.bourFlatHessianEstimate_metricFrozen, SmoothPoincare.bourUniformHarmonicMetricW2_of_flatLaplacian}
\leanok
Assume the scalar whole-space Laplacian estimate
$\|D^2v\|_{L^6}\le C_\Delta\|\Delta v\|_{L^6}$ for compactly supported smooth
$v$ and the model Haar measure. On the realized harmonic ball of
Proposition~\ref{prop:ed50-chart}, the metric entries have a common
$W^{2,6}$ bound on $B(0,5r/2)$. The constant precedes the choice of
metric, base point and realization. The statement applies to $a=1$
and $a=Q$ under the corresponding normalized ordered curvature ceiling.
\end{proposition}
\begin{proof}
\leanok
Conjugate the frozen inverse Gram operator to the Laplacian and cancel the constant
Haar Jacobian. Bound its inverse transformation by the metric upper form.
Use ellipticity and the first metric jets to choose a common freezing radius.
Absorb the oscillation term, differentiate a compact cutoff, and take a fixed finite cover.
\end{proof}

\begin{proposition}[Principal forcing for the trace-free Ricci tensor]
\label{prop:ed51-principal}
\uses{prop:ed50-chart,prop:ed50-commuted}
\lean{SmoothPoincare.bourStationary_roughTraceFreeRicci_eq, SmoothPoincare.bourChartContractedChristoffelDerivative_eq_of_laplacian, SmoothPoincare.bourChartTwo_principal_bound}
\leanok
On a constant-scalar stationary four-metric the rough trace of $Z=\mathrm{Ric}_0$
is algebraic in curvature. In harmonic coordinates on an open neighborhood,
the contracted connection derivative in the ordered Hessian of $Z$ is
controlled by raised curvature, inverse metric first derivatives, and
Christoffel symbols. The scalar principal forcing for each entry of $Z$
is bounded by a zeroth-order constant plus a multiple of the sum of its
first coordinate derivatives in absolute value. For $h=ag$ the scalar
coefficient is $(2-p)R(g)/(12a)$.
\end{proposition}
\begin{proof}
\leanok
Differentiate $h^{ab}\Gamma^k_{ab}=0$, and use the coordinate curvature
identity to switch the two derivative slots. Expand both covariant
derivatives on a two-tensor and insert the resulting identity in each
lower slot. Constant rescaling preserves the covariant tensor $Z$.
\end{proof}

\begin{proposition}[Interior absorption at exponent six]
\label{prop:ed51-cutoff}
\uses{prop:ed51-metric,prop:ed51-principal}
\lean{SmoothPoincare.Bour.sixthPowerDirectional_cutoff_lpNorm, SmoothPoincare.bourCutoffSystem_real_norm_bounds, SmoothPoincare.bourCutoffSystem_sixthInterior}
\leanok
On a single freezing patch suppose the stationary coordinate equations satisfy
$|Pu_j|\le F_0+F_1\sum_{i,k}|\partial_i u_k|$. A compact cutoff $\eta$ gives
bounds on $\|D^2(\eta^2u_j)\|_6$ and $\sum_{i,j}\|\eta\partial_i u_j\|_6$.
Their constants depend on the frozen estimate, coefficient and cutoff
bounds, patch volume and $\sup |u_j|$. The coupled system in four dimensions
has sixteen ordered tensor entries and sixty-four first derivatives.
\end{proposition}
\begin{proof}
\leanok
Integrate the sixth power of the weighted directional derivative by parts.
Young's inequality leaves an arbitrarily small multiple of the matching
$\eta^2\partial_i^2u_j$ norm. The squared-cutoff product rule bounds
these diagonal Hessians by the localized Hessians and weighted first
derivatives. Sum all directions and choose the interpolation length so
the total coefficient is at most one half.
\end{proof}

\begin{proposition}[First hits and stationarity through the final slice]
\label{prop:ed51-terminal}
\uses{prop:ed50-terminal,prop:ed50-commuted}
\lean{SmoothPoincare.Bour.originalGradient_firstHit_nativeWindows, SmoothPoincare.Bour.bourParabolicVelocity_hasDerivAt, SmoothPoincare.Bour.originalGradient_stationary_of_terminalComparison}
\leanok
Fix one parameter and a regular original flow with unbounded ordered curvature
after a regular anchor. There are cofinal first hits $(t_n,x_n)$ for any
diverging levels $Q_n$ above the anchor maximum. Each fixed normalized
backward window eventually has curvature at most one everywhere, value
one at the final base point, and vanishing integrated native gradient square.
Under time-independent compact comparison maps, a metric lower comparison,
and uniform gradient-norm comparison with a continuous limiting field, the
limiting gradient vanishes through the final slice. This transfer works
on an exhausted pointed carrier as well as a whole closed limit.
\end{proposition}
\begin{proof}
\leanok
Minimize the time projection on each compact spacetime crossing set.
The rescaled gradient has weight minus one. A point of positive limiting
gradient norm gives a compact spatial patch of positive measure and a
time interval with uniformly positive dissipation on its actual image.
Continuity also detects positivity at the right endpoint.
\end{proof}

\begin{proposition}[Outward connection forms on a finite end link]
\label{prop:ed51-pfaffian}
\uses{prop:ed50-link,prop:ed49-end}
\lean{SmoothPoincare.bourBoundaryPathTangentialCurvature_of_connectionForms, SmoothPoincare.bourImmersionNormalDefect_eq_neg_shape, SmoothPoincare.bourFrameEulerTransgression_boundary_integral, SmoothPoincare.bourEulerBulk_setIntegral_eq_scalarFlat, SmoothPoincare.bourRoundEulerShell_total}
\leanok
In an outward-adapted oriented boundary frame let $S$ be the shape
matrix and $C$ the ambient tangential curvature on complementary
wedges. The affine connection from the split to the ambient connection
gives the integrated Pfaffian coefficient $2\langle S,C\rangle+4\det S$.
The scalar-flat bulk on a finite truncation is
$E_+(K)+E_-(K)-D(K)/2$ with the physical restricted volume.
For the supplied round polar shell $0<a<b<\pi$, its two outward
boundary integrals and its bulk sum to zero.
\end{proposition}
\begin{proof}
\leanok
Metric compatibility fixes the signs of the two mixed blocks. The curvature
along the path has tangential component $C+(1-t^2)S\wedge S$ in the
complementary-wedge convention. Alternate, contract with $S$, and integrate.
Integrate the scalar-flat density on the same core.
\end{proof}

\begin{proposition}[Right slope of the positive-class value]
\label{prop:ed51-slope}
\uses{prop:ed50-second}
\lean{SmoothPoincare.Bour.bourRegularized_secant_antitone, SmoothPoincare.Bour.bourRegularizedValue_hasRightDerivative}
\leanok
After the fixed-parameter extractions, assume one sequence of actual global
regularized minima with a whole positive unit-volume Yamabe and Bach
limit and convergence of the physical chiral and ordinary Ricci integrals.
For one exact least-$D$ comparator and $J(p)=\inf_g(2(1-p)e(g)+pD(g)/2)$
on the positive metrics, $J(0)=2I$ and $J'_+(0)=D_*/2-2I$.
\end{proposition}
\begin{proof}
\leanok
Since $e(g)\ge I$, the positive right secant of each affine functional
at height $2I$ decreases with the parameter. Its infimum has the same
monotonicity. Test the upper bound on the single least-$D$ metric and
the lower bound along the returned sequence, where $D$ converges to $D_*$.
\end{proof}

\begin{proposition}[Coupled estimates on both stationary scales]
\label{prop:ed51-coupled}
\uses{prop:ed50-chart,prop:ed51-cutoff,prop:ed51-principal}
\lean{SmoothPoincare.bourStationaryRoughTraceFreeRicci_of_partialDiffeomorph_inner, SmoothPoincare.Bour.bourUniformStationaryHarmonicTraceFreeRicci_bounds, SmoothPoincare.bourStationary_harmonic_W26, SmoothPoincare.bourMixedStationary_uniformHarmonicW26}
\leanok
For a common positive Yamabe floor, scalar ceiling and normalized ordered
curvature bound, assume the Peters chart input and the scalar Euclidean
$L^6$ Laplacian estimate. One choice of radius, chart controls and finite bound
works for all stationary constant-scalar metrics with $0\le p\le1$ at both
$a=1$ and the curvature scale $a=Q\ge1$. On the same physical realization,
the sum of the $W^{2,6}$ norms of a metric component and the corresponding
covariant trace-free Ricci component is bounded on $B(0,5r/2)$.
The closed coordinate ball of radius $r$ captures the original metric
ball of radius $r/(L\sqrt a)$.
\end{proposition}
\begin{proof}
\leanok
Transport the ordered curvature and rough trace through the local isometry on
$B(0,3r)$. Bound the stationary reaction on the original carrier in an
orthonormal frame. The squared cutoff and a common oscillation radius give
uniform bounds on each smaller patch. A fixed finite cover and reconstruction
of the full first derivative give all three norms. Metric entries of any two
extensions and their jets agree on the common open physical ball.
\end{proof}

\begin{proposition}[The fixed-parameter pointed scalar vanishes]
\label{prop:ed51-scalar}
\uses{prop:ed51-terminal,prop:ed50-commuted,prop:ed48-normal}
\lean{SmoothPoincare.Bour.terminalMetric_eq_final_of_nativeSpeed, SmoothPoincare.Bour.bourCovariantGradient_trace_eq_scalarLaplacian, SmoothPoincare.bourSmoothPointedBallCutoff, SmoothPoincare.Bour.bourHarmonic_cutoff_energy, SmoothPoincare.Bour.bourPointedStationary_scalarFlat, SmoothPoincare.Bour.bourScalarFlat_stationary_zeroParameter}
\leanok
On the fixed-parameter terminal window, convergence of the metric values
on the same time-independent maps, an upper metric comparison and uniformly
vanishing native speed give a time-independent limit. A connected noncompact
boundaryless stationary four-metric with $p>0$, compact closed physical balls,
a positive unit-ball volume floor and finite $\int R^2\,dV$ has $R=0$.
At $R=0$ the same metric satisfies $G^0=0$ and
$\Delta\mathrm{Ric}+2\,\mathrm{Rm}^{\circ}(\mathrm{Ric})
=|\mathrm{Ric}|^2g/2$. On a connected closed stationary limit with
$p>0$ the scalar is constant.
\end{proposition}
\begin{proof}
\leanok
Compare two times by evaluating the actual velocity on the two transported
tangent vectors. Trace the curvature-commuted gradient to get
$-p\Delta R/4$. On the pointed carrier test harmonicity against $\eta^2 R$.
The compactly supported flux bounds $\int\eta^2|\nabla R|^2$ by
$4\sup|\nabla\eta|^2\int R^2$. Smooth distance cutoffs at increasing
radii and Fatou give $\nabla R=0$. Disjoint unit balls force infinite
volume, so the square-integrable constant vanishes.
\end{proof}


\begin{proposition}[Interior control of the first derivatives]
\label{prop:ed51-morrey}
\uses{prop:ed51-coupled}
\lean{SmoothPoincare.bourW26Norm_volume_eq, SmoothPoincare.bourHarmonicMorrey_of_W26, SmoothPoincare.bourMixedStationary_uniformHarmonicC1Third}
\leanok
In the four-dimensional Euclidean model, on $B(0,2r)$ each metric entry
and each covariant trace-free Ricci entry of the same physical realization
has uniformly bounded value and first derivative. Both first derivatives
are uniformly H\"older continuous with exponent $1/3$. The point, metric,
scale and regularization parameter are selected after the common bound.
\end{proposition}
\begin{proof}
\leanok
Normalize the additive Haar measure and localize the $W^{2,6}$ bound on
$B(0,5r/2)$ using a smooth cutoff equal to one on the smaller ball.
Apply the Euclidean Morrey inequalities to the entry and its first
derivative, with $1-4/6=1/3$.
\end{proof}
